\documentclass[11pt]{article}
\usepackage[margin=1in]{geometry}
\usepackage{enumitem}
% =============================================================================
% Preambles/header.tex — shared LaTeX/Beamer preamble for this project
%
% Use from a Beamer lecture by prepending:
%     \input{header}
% and compiling with TEXINPUTS=../Preambles:$TEXINPUTS (the /compile-latex
% skill does this automatically).
%
% PALETTE CONTRACT. Color names defined here MUST match the names in
% Quarto/theme-template.scss so Beamer and Quarto renderings use the same
% palette. The scripts/check-palette-sync.sh script enforces this.
%
% To customize: edit the HEX values below AND the matching $variable in
% Quarto/theme-template.scss, then run scripts/check-palette-sync.sh.
% =============================================================================

% -----------------------------------------------------------------------------
% Engine detection + encoding
%
% Our /compile-latex skill uses XeLaTeX, where inputenc is unnecessary and
% can produce encoding warnings. Only load inputenc under pdfTeX.
% -----------------------------------------------------------------------------
\usepackage{iftex}
\ifPDFTeX
  \usepackage[utf8]{inputenc}
\fi

\usepackage{amsmath, amssymb, amsthm}
\usepackage{booktabs}
\usepackage{graphicx}

% -----------------------------------------------------------------------------
% PALETTE — mirrors Quarto/theme-template.scss variable names.
% Load xcolor and define colors BEFORE anything that references them
% (notably hyperref, hypersetup, and the Beamer color assignments).
% -----------------------------------------------------------------------------
\usepackage{xcolor}

\definecolor{primary-blue}{HTML}{012169}      % headings, accents
\definecolor{primary-gold}{HTML}{B9975B}      % emphasis, borders
\definecolor{highlight-yellow}{HTML}{F2A900}  % markers, alerts
\definecolor{light-bg}{HTML}{E8EDF5}          % title / section backgrounds
\definecolor{jet}{HTML}{1A1A1A}               % body text

% Semantic colors (used in diagrams and annotations)
\definecolor{positive}{HTML}{15803D}          % good / observed / identified
\definecolor{negative}{HTML}{B91C1C}          % bad / problematic / bias
\definecolor{neutral}{HTML}{525252}           % reference / context

% Highlight accents (match .hi-* classes in the SCSS)
\definecolor{hi-slate}{HTML}{314F4F}
\definecolor{hi-green}{HTML}{2E8B57}
\definecolor{hi-red}{HTML}{C41E3A}

% Semantic aliases so skill templates and snippets can write
% \textcolor{accent}{...} without hard-coding "primary-blue".
\colorlet{accent}{primary-blue}
\colorlet{accent2}{primary-gold}

% -----------------------------------------------------------------------------
% Hyperref — loaded AFTER xcolor + palette so link colors resolve.
% -----------------------------------------------------------------------------
\usepackage{hyperref}
\hypersetup{colorlinks=true, linkcolor=primary-blue, urlcolor=primary-blue,
            citecolor=primary-blue, pdfborder={0 0 0}}

% -----------------------------------------------------------------------------
% Course code — set once per deck (after \input{header}, before \begin{document})
% via \coursecode{CS401}. Shown in the footer on every slide so a deck's course
% is identifiable without opening the title page. Empty by default.
% -----------------------------------------------------------------------------
\makeatletter
\newcommand{\coursecode}[1]{\gdef\@coursecodetext{#1}}
\gdef\@coursecodetext{}
\makeatother

% -----------------------------------------------------------------------------
% Beamer theme assignments (only apply under Beamer).
%
% We detect Beamer via \@ifclassloaded{beamer}, which is the documented,
% non-internal way. This must be wrapped in \makeatletter/\makeatother
% because \@ifclassloaded contains an @-catcode character.
% -----------------------------------------------------------------------------
\makeatletter
\@ifclassloaded{beamer}{%
  \usefonttheme{professionalfonts}
  \setbeamercolor{structure}{fg=primary-blue}
  \setbeamercolor{frametitle}{fg=primary-blue}
  \setbeamercolor{title}{fg=primary-blue}
  \setbeamercolor{subtitle}{fg=primary-gold}
  \setbeamercolor{author}{fg=primary-blue}
  \setbeamercolor{institute}{fg=neutral}
  \setbeamercolor{date}{fg=primary-gold}
  \setbeamercolor{itemize item}{fg=primary-blue}
  \setbeamercolor{itemize subitem}{fg=primary-gold}
  \setbeamercolor{alerted text}{fg=primary-gold}
  \setbeamercolor{block title}{fg=white, bg=primary-blue}
  \setbeamercolor{block body}{bg=light-bg}
  % Semantic block variants: block=definition/concept (blue), exampleblock=
  % worked example (green/positive), alertblock=key takeaway or caution (gold).
  % Keep to <=2 colored boxes per slide across ALL three variants (INV-7).
  \setbeamercolor{block title example}{fg=white, bg=positive}
  \setbeamercolor{block body example}{bg=light-bg}
  \setbeamercolor{block title alerted}{fg=white, bg=primary-gold}
  \setbeamercolor{block body alerted}{bg=light-bg}

  % Minimal footer: course code (left) + page X / N (right), no navigation clutter
  \setbeamertemplate{navigation symbols}{}
  \setbeamertemplate{footline}{%
    \color{neutral}\footnotesize\hspace{0.7em}\@coursecodetext\hfill
    \insertframenumber/\inserttotalframenumber\hspace{0.7em}\vspace{0.3em}}
}{}
\makeatother

% -----------------------------------------------------------------------------
% TikZ setup — libraries the snippet gallery uses
% -----------------------------------------------------------------------------
\usepackage{tikz}
\usetikzlibrary{arrows.meta, positioning, calc, decorations.pathreplacing,
                fit, shapes.geometric, backgrounds}

% Shared TikZ styles — snippets and hand-written diagrams can pick these up
% via `\begin{tikzpicture}[every node/.style={...}]` or by naming specific
% styles (e.g., `\node[dag-node] (X) at (0,0) {X};`).
\tikzset{
  % Node styles for causal DAGs and similar
  dag-node/.style = {
    draw=primary-blue, fill=light-bg, rounded corners=2pt,
    minimum width=1.2cm, minimum height=0.9cm, align=center, font=\small
  },
  decision-node/.style = {
    draw=primary-gold, fill=white, diamond, aspect=1.8,
    minimum width=1.8cm, minimum height=1.2cm, align=center, font=\small,
    inner sep=1pt
  },
  % Edge styles — observed vs counterfactual
  observed-edge/.style = {-{Latex[length=2.5mm]}, thick, draw=jet},
  counterfactual-edge/.style = {-{Latex[length=2.5mm]}, thick, draw=neutral, dashed},
  confound-edge/.style = {-{Latex[length=2.5mm]}, thick, draw=negative, dashed},
  % Data-point styles for plots
  observed-dot/.style = {fill=primary-blue, draw=primary-blue, circle, inner sep=1.2pt},
  counterfactual-dot/.style = {fill=white, draw=neutral, circle, inner sep=1.2pt}
}

% -----------------------------------------------------------------------------
% Convenience macros
% -----------------------------------------------------------------------------
\newcommand{\muted}[1]{\textcolor{neutral}{#1}}
\newcommand{\key}[1]{\textcolor{primary-gold}{\textbf{#1}}}
\newcommand{\good}[1]{\textcolor{positive}{#1}}
\newcommand{\bad}[1]{\textcolor{negative}{#1}}

% Standout frame macro: full-bleed dark background for section transitions.
% Beamer-only (uses \begin{frame}/\setbeamercolor/\usebeamerfont) -- do not
% call from an article-class file (e.g. Notes/<CODE>/*-notes.tex). \muted,
% \key, \good, \bad, and \coursecode above are plain xcolor/gdef and are
% safe to use from either class; verified when Notes/article-class support
% was added.
% Expand: \transitionslide{Your transition title}
\newcommand{\transitionslide}[1]{%
  \begingroup
  \setbeamercolor{background canvas}{bg=primary-blue}
  \begin{frame}[plain]
    \centering
    \vfill
    {\usebeamerfont{title}\color{white}\Large #1\par}
    \vfill
  \end{frame}
  \endgroup
}

% Section divider frame (Beamer-only), an alternative to \transitionslide with a
% darker two-line style: near-black (jet) background, a small white label line
% above a larger gold title, and a thin gold rule along the bottom edge.
% Expand: \sectiondivider{Section 2}{Pointers}
\newcommand{\sectiondivider}[2]{%
  \begingroup
  \setbeamercolor{background canvas}{bg=jet}
  \begin{frame}[plain]
    \vspace*{3.0cm}
    {\color{white}\ttfamily\large #1\par}
    \vspace{0.5cm}
    {\color{primary-gold}\LARGE #2\par}
    \begin{tikzpicture}[remember picture, overlay]
      \draw[primary-gold, line width=2.5pt]
        ($(current page.south west)+(0,0.1)$) -- ($(current page.south east)+(0,0.1)$);
    \end{tikzpicture}
  \end{frame}
  \endgroup
}

% -----------------------------------------------------------------------------
% Channel watermark — "Concepts That Click" (YouTube).
%
% Article-class files (CompetitiveExam Q/A sets, Notes, Assignments) get a
% light diagonal draft watermark on every page plus a centered footer naming
% the channel. Beamer decks get a subtle TikZ background watermark instead
% (the footline already carries the course code). Centralized here so every
% MCQ PDF is branded without per-file edits.
% -----------------------------------------------------------------------------
\makeatletter
\@ifclassloaded{article}{%
  \usepackage{draftwatermark}
  \SetWatermarkText{Concepts That Click}
  \SetWatermarkScale{1}
  \SetWatermarkFontSize{2cm}
  \SetWatermarkLightness{0.86}
  \SetWatermarkAngle{45}
  \usepackage{fancyhdr}
  \pagestyle{fancy}
  \fancyhf{}
  \fancyfoot[C]{\footnotesize\textcolor{neutral}{Concepts That Click~|~YouTube: \href{https://www.youtube.com/@concepts.thatclick}{@concepts.thatclick}}}
  \renewcommand{\headrulewidth}{0pt}
  \renewcommand{\footrulewidth}{0pt}
  \fancypagestyle{plain}{%
    \fancyhf{}%
    \fancyfoot[C]{\footnotesize\textcolor{neutral}{Concepts That Click~|~YouTube: \href{https://www.youtube.com/@concepts.thatclick}{@concepts.thatclick}}}%
    \renewcommand{\headrulewidth}{0pt}%
    \renewcommand{\footrulewidth}{0pt}%
  }
}{}
\@ifclassloaded{beamer}{%
  \setbeamertemplate{background}{%
    \begin{tikzpicture}[remember picture, overlay]
      \node[rotate=45, opacity=0.055, align=center]
        at (current page.center)
        {\Huge\textcolor{neutral}{Concepts That Click}};
    \end{tikzpicture}%
  }
}{}
\makeatother 
\coursecode{JKSSB}
\title{Lesson 64 Practice: Current Cybersecurity, Cryptography and Cyber Law}
\author{JKSSB Computer Awareness}
\date{\today}
\begin{document}
\maketitle
\noindent\textbf{Provenance:} All 20 questions are original JKSSB-pattern
questions (\textit{[Original]}); none reproduces an official paper item.

\begin{enumerate}[label=\textbf{Q\arabic*.}, leftmargin=*]
\item \textit{[Original]} A parasitic (file-infector) virus is best described as:
  \begin{enumerate}[label=\Alph*.]
    \item a virus that attaches itself to an executable file and runs when the host program runs
    \item a worm that spreads across a network without any host program
    \item a Trojan horse that disguises itself as useful software without self-replicating
    \item a flood of traffic from many machines that overwhelms a server
  \end{enumerate}
\item \textit{[Original]} Expand DES.
  \begin{enumerate}[label=\Alph*.]
    \item Data Encryption Standard
    \item Digital Encryption System
    \item Distributed Encoding Scheme
    \item Data Evaluation Service
  \end{enumerate}
\item \textit{[Original]} DES processes data in blocks of:
  \begin{enumerate}[label=\Alph*.]
    \item 32 bits
    \item 56 bits
    \item 64 bits
    \item 128 bits
  \end{enumerate}
\item \textit{[Original]} Which of the following is \textbf{NOT} true of a digital signature?
  \begin{enumerate}[label=\Alph*.]
    \item It proves the sender's authenticity.
    \item It proves the message was not altered in transit.
    \item It provides integrity of the message.
    \item It encrypts the message body so only the recipient can read it.
  \end{enumerate}
\item \textit{[Original]} Encryption primarily provides \underline{\hspace{2.2cm}}, whereas a digital signature primarily provides \underline{\hspace{2.2cm}}.
  \begin{enumerate}[label=\Alph*.]
    \item confidentiality; authenticity and integrity
    \item authenticity; confidentiality
    \item availability; confidentiality
    \item integrity; availability
  \end{enumerate}
\item \textit{[Original]} Which firewall type examines only per-packet headers such as source/destination address, port, and protocol?
  \begin{enumerate}[label=\Alph*.]
    \item Packet-filtering firewall
    \item Stateful-inspection firewall
    \item Application (proxy) firewall
    \item Personal anti-virus program
  \end{enumerate}
\item \textit{[Original]} A firewall that remembers open sessions and drops packets that do not belong to a known session is a:
  \begin{enumerate}[label=\Alph*.]
    \item packet-filtering firewall
    \item stateful-inspection firewall
    \item boot-sector virus
    \item Trojan horse
  \end{enumerate}
\item \textit{[Original]} Which firewall acts as a go-between, inspecting application-layer content and commands?
  \begin{enumerate}[label=\Alph*.]
    \item Packet-filtering firewall
    \item Stateful-inspection firewall
    \item Application (proxy) firewall
    \item Spam filter of the mail client only
  \end{enumerate}
\item \textit{[Original]} A fraudulent mail styled as a bank KYC notice asks a clerk for her OTP on a fake login page. This is:
  \begin{enumerate}[label=\Alph*.]
    \item phishing
    \item ransomware
    \item DDoS
    \item disk defragmentation
  \end{enumerate}
\item \textit{[Original]} A clerk's files are suddenly unreadable, with a message demanding payment for the decryption key. This is:
  \begin{enumerate}[label=\Alph*.]
    \item phishing
    \item ransomware
    \item DDoS
    \item hacking without any malicious goal
  \end{enumerate}
\item \textit{[Original]} Which of the following correctly distinguishes the four threats?
  \begin{enumerate}[label=\Alph*.]
    \item Phishing floods a server; ransomware tricks the user; DDoS encrypts files; hacking defragments disks.
    \item Phishing tricks the user into revealing credentials; ransomware locks files for ransom; DDoS floods a service; hacking is unauthorised access generally.
    \item Phishing encrypts files; ransomware floods a server; DDoS reveals OTPs; hacking is a firewall type.
    \item Phishing is a firewall; ransomware is an anti-virus; DDoS is encryption; hacking is a backup.
  \end{enumerate}
\item \textit{[Original]} Expand DDoS.
  \begin{enumerate}[label=\Alph*.]
    \item Distributed Denial of Service
    \item Direct Disk Operating System
    \item Digital Data on Server
    \item Dynamic Domain of Service
  \end{enumerate}
\item \textit{[Original]} Evaluate the statements:
  \begin{enumerate}[label=\Roman*.]
    \item DES is a symmetric cipher: the same key encrypts and decrypts.
    \item DES works on 64-bit blocks of data.
    \item A digital signature by itself keeps the message body secret from everyone except the recipient.
  \end{enumerate}
  \begin{enumerate}[label=\Alph*.]
    \item I only
    \item I and II only
    \item II and III only
    \item I, II, and III
  \end{enumerate}
\item \textit{[Original]} Which pairing of firewall type and function is correctly matched?
  \begin{enumerate}[label=\Alph*.]
    \item Packet-filtering --- inspects application content as a go-between
    \item Stateful inspection --- tracks connection state in addition to headers
    \item Application (proxy) --- decides each packet by headers only
    \item Packet-filtering --- disinfects infected executable files
  \end{enumerate}
\item \textit{[Original]} Which of the following is a governance concern of the Internet at awareness level?
  \begin{enumerate}[label=\Alph*.]
    \item Which country's law applies when an act crosses borders
    \item The colour scheme of a single office website
    \item The brand of printer attached to the office desktop
    \item The seating plan of the computer laboratory
  \end{enumerate}
\item \textit{[Original]} Which data-protection principle requires that personal data be used only for the purpose stated at collection?
  \begin{enumerate}[label=\Alph*.]
    \item Purpose limitation (stated-purpose use)
    \item Unlimited retention
    \item Open sharing without consent
    \item Default publication
  \end{enumerate}
\item \textit{[Original]} Copying licensed software without permission is, at awareness level, an example of:
  \begin{enumerate}[label=\Alph*.]
    \item copyright infringement
    \item fair dealing for private study
    \item first ownership
    \item lawful backup under licence
  \end{enumerate}
\item \textit{[Original]} Match the copyright concept to what it covers: (a) meaning of copyright, (b) first ownership, (c) infringement, (d) fair dealing. Which match is correct?
  \begin{enumerate}[label=\Alph*.]
    \item (a)---punishment for wilful violation; (b)---permitted study uses; (c)---who owns first; (d)---owner's exclusive rights
    \item (a)---owner's exclusive rights; (b)---who owns the work first; (c)---what counts as violation; (d)---limited permitted uses such as study
    \item (a)---who owns first; (b)---owner's rights; (c)---permitted uses; (d)---what counts as violation
    \item (a)---permitted uses; (b)---violation; (c)---owner's rights; (d)---first ownership
  \end{enumerate}
\item \textit{[Original]} Expand ICANN.
  \begin{enumerate}[label=\Alph*.]
    \item Internet Corporation for Assigned Names and Numbers
    \item International Computer and Network Node
    \item Internet Code for Access and Navigation
    \item Indian Council for Assigned Network Names
  \end{enumerate}
\item \textit{[Original]} An office server slows to a halt under a flood of traffic from many machines, so genuine users cannot reach it. The correct response pair is:
  \begin{enumerate}[label=\Alph*.]
    \item identify it as DDoS and filter/block the flood at the network barrier (firewall/ISP level)
    \item identify it as phishing and ask users to change their OTPs only
    \item identify it as a file-infector virus and reinstall the printer driver
    \item identify it as fair dealing and take no technical action
  \end{enumerate}
\end{enumerate}
\end{document}
